\section{Introduction}
Image classifiers often learn shortcuts that work on most training images and fail on a minority of them \cite{geirhos2020shortcut}. Waterbirds \cite{sagawa2020groupdro} makes this concrete: birds from CUB \cite{wah2011cub} are pasted onto Places backgrounds \cite{zhou2018places} so that 95\% of training waterbirds appear on water and 95\% of landbirds on land. The standard metric is worst-group accuracy (WGA), the lowest accuracy over the four bird-by-background groups. The two smallest test groups contain 642 images each, so one image moves WGA by 0.16 percentage points (pp), and recent methods differ by a few tenths of a point.

Much of the recent progress comes from methods that keep an ERM backbone frozen and fit a light head on the small validation set with group labels \cite{kirichenko2023dfr,izmailov2022feature,labonte2023fewer}. Diverse prototypical ensembles (DPE) \cite{to2025dpe} are a strong example: fifteen prototype heads are trained in sequence on group-balanced subsets of the 1,199 annotated validation images, and their probabilities are averaged. Because the validation set is the whole of the supervision, it is natural to try to use it more carefully after the heads are trained: to learn how to combine the heads \cite{wolpert1992stacking,breiman1996stacked}, or to recalibrate the ensemble on out-of-fold scores obtained by cross-fitting \cite{chernozhukov2018dml,platt1999probabilistic}. We set out to measure whether such post-hoc procedures improve DPE when every choice is made on validation data only.

On the official test labels, they do not. With validation-selected settings, cross-fitted recalibration, a learned combiner and a recalibration map transferred to the full ensemble are all 0.2--0.9\pp{} below DPE fitted on all images. We traced this to a known label error \cite{sagawa2020code,asgari2022masktune}. The generating script marks a CUB species as a waterbird if its name contains a keyword such as ``tern'', which also occurs in ``Eastern'' and ``Western'', so Eastern Towhee, Western Meadowlark and Western Wood-Pewee, which are landbirds, are labelled waterbird: 18 validation and 90 test images. Removing the 90 test images, or relabelling them, changes the conclusion: on the 5,704 correctly labelled test images, the same predictions beat DPE by 1.1--1.7\pp{} on all five standard-ERM backbones. The reason is that the error falls on the groups that define WGA. The 43 mislabelled waterbird-on-land test images make that group look like DPE's worst group; on correctly labelled images the worst group is landbird-on-water, which is the group the post-hoc procedures improve.

Our contributions are:
\begin{itemize}
\item \textbf{A controlled evaluation of post-hoc procedures for DPE}: settings chosen on validation data only, strictly nested inner cross-validation in which every head and calibrator is refitted without the images it is scored on, ten split seeds, and paired per-backbone statistics. Every prediction is scored on official labels, corrected labels and the correctly labelled subset.
\item \textbf{Post-hoc procedures improve DPE's real worst group.} On correctly labelled test images, a linear combiner of head scores selected on validation data gains $+1.74$\pp{} ($[+1.45,+2.03]$), cross-fitted recalibration $+1.11$\pp{} ($[+0.67,+1.56]$) and a map learnt on out-of-fold scores and applied to full-validation DPE $+1.39$\pp{} ($[+1.13,+1.65]$), each on all five standard-ERM backbones. The combiner also beats scalar recalibration of the same heads ($+1.16$\pp), so head disagreements carry information beyond a shift of the ensemble score.
\item \textbf{Official test labels hide these gains.} The 90 mislabelled test images misidentify DPE's worst group, and official-label results score the improvements as losses. To our knowledge, previous discussions of this error \cite{sagawa2020code,asgari2022masktune} do not separate its effect through the test labels from its effect through the validation labels on which methods such as DPE are fitted.
\item \textbf{Validation labels are a separate lever.} Refitting DPE after correcting the 18 validation images gains $+2.30$\pp{} on correctly labelled test images ($[+1.65,+2.94]$), and $+2.82$\pp{} on four fresh backbones trained with the authors' code, with the contrast declared before that run.
\end{itemize}
Not every procedure helps: holding images back without a combiner, uniform cross-fitting and an exact WGA threshold do not. We also reproduce DPE with the authors' code on two fresh ERM backbones (89.4\% and 86.5\% WGA against 91.0\% reported).

\textbf{Scope.} We do not propose a new method or claim a state of the art, and the study covers one dataset and one method family. The work went through three stages that differ in how independent they are of the test set: exploratory development on seven saved backbones, during which the test set was inspected; a strictly nested rerun in which all settings were chosen on validation data; and a replication of the label effect on fresh backbones. The scoring of the nested rerun on correctly labelled images was done after that run, without refitting or reselecting anything. Code, per-run results and audits are released.\footnote{\url{https://github.com/chautanphat/waterbirds-shortcut-mitigation}}
